\documentclass{article}
\usepackage[T1]{fontenc}
\usepackage{lmodern}
\usepackage{graphicx}
\usepackage{amsmath,amssymb}
\usepackage[a4paper,margin=2cm]{geometry}
\usepackage[absolute,overlay]{textpos}
\setlength{\TPHorizModule}{1mm}
\setlength{\TPVertModule}{1mm}
\usepackage[indonesian]{babel}
\usepackage{hyperref}
\setlength{\emergencystretch}{2em}
% unit: Halaman 24

\begin{document}

% === Halaman 24 (llm) ===

\section*{2. Plaintext-Recovery Attack}

\begin{itemize}
    \item Tujuan: memperoleh plainteks, tetapi tidak harus memperoleh kunci.
    \item Contoh: Misalkan sebuah cipherteks dikirimir:
    \[ C = XQZ\dots \]
    Penyerang menganalisis struktur cipherteks dan mengetahui bahwa pesan tersebut kemungkinan berbunyi:
    \[ \text{MEET AT THE STATION} \]
    tanpa mengetahui kunci sebenarnya.
\end{itemize}

Dalam beberapa situasi, mendapatkan plainteks tertentu sudah cukup bagi penyerang. Misalnya, jika pesan berisi informasi penting, penyerang mungkin tidak perlu mengetahui kunci untuk mencapai tujuannya.

\end{document}
